\documentclass[border=12pt]{standalone}
\usepackage[no-math]{fontspec}
\usepackage[brazil]{babel}
\usepackage{amsmath,amssymb}
\IfFileExists{C:/Windows/Fonts/Lato_Regular.ttf}{%
  \setmainfont{Lato}[
    Path = C:/Windows/Fonts/,
    UprightFont = Lato_Regular.ttf,
    BoldFont = Lato_Bold.ttf,
    ItalicFont = Lato_Italic.ttf,
    BoldItalicFont = Lato_BoldItalic.ttf
  ]
  \setsansfont{Lato}[
    Path = C:/Windows/Fonts/,
    UprightFont = Lato_Regular.ttf,
    BoldFont = Lato_Bold.ttf,
    ItalicFont = Lato_Italic.ttf,
    BoldItalicFont = Lato_BoldItalic.ttf
  ]
}{%
  \RequirePackage[default]{lato}
}
\renewcommand{\familydefault}{\sfdefault}

\usepackage{xcolor}
\usepackage{tikz}
\usepackage{pgfplots}
\pgfplotsset{compat=1.18}
\usetikzlibrary{arrows.meta,positioning,shapes.geometric,fit,backgrounds,calc,decorations.pathreplacing,patterns}

% Paleta institucional UEFS / Sociologia Rural
\definecolor{srNavy}{HTML}{1F4268}
\definecolor{srNavyDark}{HTML}{152E4A}
\definecolor{srGold}{HTML}{B38A3D}
\definecolor{srGoldDark}{HTML}{8C6927}
\definecolor{srTeal}{HTML}{0E7C7B}
\definecolor{srForest}{HTML}{2D6A4F}
\definecolor{srAlert}{HTML}{9E2A2B}
\definecolor{srInk}{HTML}{26313D}
\definecolor{srMuted}{HTML}{5B6573}

\definecolor{bgNavy}{HTML}{EDF2F7}
\definecolor{bgGold}{HTML}{F8F3E9}
\definecolor{bgTeal}{HTML}{EBF6F5}
\definecolor{bgForest}{HTML}{EEF5F0}
\definecolor{bgAlert}{HTML}{FDF2F2}

\tikzset{
  >=Stealth,
  srBox/.style={
    draw=#1, fill=white, rounded corners=5pt, line width=1.15pt,
    inner sep=7pt, align=center, font=\footnotesize\sffamily, text=srInk
  },
  srTint/.style={
    draw=#1, fill=#1!8!white, rounded corners=5pt, line width=1.15pt,
    inner sep=7pt, align=center, font=\footnotesize\sffamily, text=srInk
  },
  srBadge/.style={
    fill=#1, text=white, rounded corners=9pt, inner xsep=8pt, inner ysep=3pt,
    font=\scriptsize\bfseries\sffamily
  },
  srTitle/.style={font=\normalsize\bfseries\sffamily, text=srNavy, align=center},
  srSub/.style={font=\footnotesize\sffamily, text=srMuted, align=center},
  srArrow/.style={->, >=Stealth, line width=1.35pt},
  srDArrow/.style={<->, >=Stealth, line width=1.35pt}
}


\begin{document}
\begin{tikzpicture}[
  card/.style={
    draw=#1,
    fill=white,
    rounded corners=6pt,
    line width=1.2pt,
    minimum width=5.55cm,
    minimum height=5.15cm,
    inner sep=0pt
  },
  cardHead/.style={
    fill=#1,
    text=white,
    font=\small\bfseries\sffamily,
    align=center,
    minimum width=5.55cm,
    minimum height=0.95cm,
    rounded corners=5pt
  },
  cardBody/.style={
    anchor=north,
    text width=4.85cm,
    align=left,
    font=\footnotesize\sffamily,
    text=srInk,
    inner sep=0pt
  }
]
  % Título principal
  \node[srTitle, font=\large\bfseries\sffamily] (title) at (0, 4.10)
    {Diversidade Étnico-Racial e Territórios Tradicionais no Brasil (Decreto nº 6.040/2007)};

  % Posições dos 3 cartões
  \coordinate (c1) at (-5.95, 0.85);
  \coordinate (c2) at ( 0.00, 0.85);
  \coordinate (c3) at ( 5.95, 0.85);

  % --- Cartão 1: Povos Indígenas ---
  \node[card=srForest, fill=bgForest!55!white] (box1) at (c1) {};
  \node[cardHead=srForest, anchor=north] (head1) at (box1.north) {};
  \fill[srForest] ($(head1.south west)+(0,0.18cm)$) rectangle (head1.south east);
  \node[font=\small\bfseries\sffamily, text=white] at (head1) {Povos Indígenas};
  \node[cardBody] at ($(head1.south)+(0,-0.32cm)$) {%
    \hyphenpenalty=10000\exhyphenpenalty=10000
    \textbf{\textcolor{srForest}{\textbullet}} \textbf{Art. 231 da CF/1988}: direitos originários sobre as terras que tradicionalmente ocupam.\\[7pt]
    \textbf{\textcolor{srForest}{\textbullet}} \textbf{Terras Indígenas demarcadas}: proteção territorial e usufruto exclusivo.\\[7pt]
    \textbf{\textcolor{srForest}{\textbullet}} \textbf{Sociobiodiversidade}: sistemas tradicionais e manejo da biodiversidade.
  };

  % --- Cartão 2: Comunidades Quilombolas ---
  \node[card=srGoldDark, fill=bgGold!70!white] (box2) at (c2) {};
  \node[cardHead=srGoldDark, anchor=north] (head2) at (box2.north) {};
  \fill[srGoldDark] ($(head2.south west)+(0,0.18cm)$) rectangle (head2.south east);
  \node[font=\small\bfseries\sffamily, text=white] at (head2) {Comunidades Quilombolas};
  \node[cardBody] at ($(head2.south)+(0,-0.32cm)$) {%
    \hyphenpenalty=10000\exhyphenpenalty=10000
    \textbf{\textcolor{srGoldDark}{\textbullet}} \textbf{Art. 68 do ADCT} \& \textbf{Decreto nº 4.887/2003}: reconhecimento e autodefinição.\\[7pt]
    \textbf{\textcolor{srGoldDark}{\textbullet}} \textbf{Titulação coletiva (INCRA)}: propriedade definitiva em nome da associação.\\[7pt]
    \textbf{\textcolor{srGoldDark}{\textbullet}} \textbf{Território e cultura}: memória territorial e roças tradicionais.
  };

  % --- Cartão 3: Povos e Comunidades Tradicionais ---
  \node[card=srTeal, fill=bgTeal!60!white] (box3) at (c3) {};
  \node[cardHead=srTeal, anchor=north] (head3) at (box3.north) {};
  \fill[srTeal] ($(head3.south west)+(0,0.18cm)$) rectangle (head3.south east);
  \node[font=\small\bfseries\sffamily, text=white] at (head3) {Povos e Comunidades Tradicionais};
  \node[cardBody] at ($(head3.south)+(0,-0.32cm)$) {%
    \hyphenpenalty=10000\exhyphenpenalty=10000
    \textbf{\textcolor{srTeal}{\textbullet}} \textbf{Fundos e Fechos de Pasto}: uso comum da terra no semiárido e cerrado.\\[7pt]
    \textbf{\textcolor{srTeal}{\textbullet}} \textbf{Catadoras de Mangaba}: extrativismo costeiro e soberania alimentar.\\[7pt]
    \textbf{\textcolor{srTeal}{\textbullet}} \textbf{Ribeirinhos, Extrativistas e Geraizeiros}: saberes ecológicos locais.
  };

  % Banner inferior de síntese
  \node[
    draw=srNavy,
    fill=bgNavy,
    rounded corners=6pt,
    line width=1.25pt,
    minimum width=17.45cm,
    minimum height=1.05cm,
    align=center,
    font=\small\sffamily,
    text=srNavyDark,
    inner xsep=12pt,
    inner ysep=8pt
  ] (banner) at (0, -2.90) {%
    \textit{\textbf{PNPCT (Decreto nº 6.040/2007):} proteção territorial, autodeterminação e coexistência de modos de vida rurais.}
  };

  % Conectores convergentes para o banner de síntese
  \draw[srArrow, srForest] (box1.south) -- (box1.south |- banner.north);
  \draw[srArrow, srGoldDark] (box2.south) -- (box2.south |- banner.north);
  \draw[srArrow, srTeal] (box3.south) -- (box3.south |- banner.north);
\end{tikzpicture}
\end{document}
